
\chapter*{专有术语英汉对照表}

\addcontentsline{toc}{chapter}{专有术语英汉对照表}
\markboth{专有术语英汉对照表}{专有术语英汉对照表}

本书翻译严格遵循统一的理论核物理中文学术术语体系。


\section*{1. 核心模型、方法与理论框架 (Core Theories \& Frameworks)}

\begin{longtable}{p{4.2cm} p{4.2cm} p{1.6cm} p{4.5cm}}
\toprule
\textbf{英文原名 (En)} & \textbf{规范中文译名 (Zh)} & \textbf{缩写/记号} & \textbf{重点说明 / 典型章节} \tabularnewline
\midrule
\endfirsthead
\toprule
\textbf{英文原名 (En)} & \textbf{规范中文译名 (Zh)} & \textbf{缩写/记号} & \textbf{重点说明 / 典型章节} \tabularnewline
\midrule
\endhead
\bottomrule
\endfoot
\bottomrule
\endlastfoot
Many-Body Problem & 多体问题 &  & 全书核心主题 \tabularnewline
Liquid Drop Model & 液滴模型 & LDM & 第 1 章 \tabularnewline
Shell Model & 壳模型 & SM & 第 2 章 \tabularnewline
Deformed Shell Model & 形变壳模型 &  & 第 2 章，Nilsson 势等 \tabularnewline
Mean Field Approximation & 平均场近似 &  & 贯穿全书 \tabularnewline
Self-Consistent Field & 自洽场 & SCF & 第 5、7 章 \tabularnewline
Hartree-Fock Method & Hartree-Fock 方法 / HF 方法 & HF & 第 5 章 \tabularnewline
Density-Dependent Hartree-Fock & 密度依赖的 Hartree-Fock 方法 & DDHF & 第 5 章，Skyrme / Gogny 力 \tabularnewline
Local Density Approximation & 局域密度近似 & LDA & 第 5、13 章 \tabularnewline
Pairing Correlation & 对关联 &  & 第 6 章 \tabularnewline
Superfluid Nuclei & 超流核 &  & 第 6 章 \tabularnewline
Seniority Model & 辛弱数模型 &  & 第 6 章，准自旋方法 \tabularnewline
BCS Theory & BCS 理论 & BCS & Bardeen-Cooper-Schrieffer 理论，第 6 章 \tabularnewline
Generalized Single-Particle Model & 推广的单粒子模型 &  & 第 7 章 \tabularnewline
Hartree-Fock-Bogoliubov Theory & Hartree-Fock-Bogoliubov 理论 & HFB & 第 7 章 \tabularnewline
Constrained Hartree-Fock / HFB & 约束 Hartree-Fock / 约束 HFB 理论 & CHF / CHFB & 第 7 章 \tabularnewline
Harmonic Vibrations & 简谐振动 &  & 第 8 章 \tabularnewline
Random Phase Approximation & 随机相位近似 & RPA & 第 8 章 \tabularnewline
Quasiparticle RPA & 准粒子随机相位近似 & QRPA & 第 8 章 \tabularnewline
Boson Expansion Method & 玻色展开方法 & BEM & 第 9 章 \tabularnewline
Generator Coordinate Method & 生成坐标法 & GCM & 第 10 章 \tabularnewline
Restoration of Broken Symmetries & 对称性的恢复 &  & 第 11 章 \tabularnewline
Spontaneous Symmetry Breaking & 自发对称性破缺 & SSB & 第 11 章 \tabularnewline
Projection Method & 投影方法 &  & 第 11 章，粒子数投影、角动量投影 \tabularnewline
Time-Dependent Hartree-Fock & 时间相关 Hartree-Fock 方法 & TDHF & 第 12 章 \tabularnewline
Semiclassical Methods & 半经典方法 &  & 第 13 章 \tabularnewline
\end{longtable}

\section*{2. 微观多体物理与算符代数 (Many-Body Physics \& Second Quantization)}

\begin{longtable}{p{4.8cm} p{4.8cm} p{5.0cm}}
\toprule
\textbf{英文原名 (En)} & \textbf{规范中文译名 (Zh)} & \textbf{常见符号 / 备注} \tabularnewline
\midrule
\endfirsthead
\toprule
\textbf{英文原名 (En)} & \textbf{规范中文译名 (Zh)} & \textbf{常见符号 / 备注} \tabularnewline
\midrule
\endhead
\bottomrule
\endfoot
\bottomrule
\endlastfoot
Second Quantization & 二次量子化 & 附录 C \tabularnewline
Creation Operator & 产生算符 & 粒子 $c_k^\dagger, a_k^\dagger$；准粒子 $\alpha_k^\dagger, \beta_k^\dagger$ \tabularnewline
Annihilation Operator & 湮灭算符 & 粒子 $c_k, a_k$；准粒子 $\alpha_k, \beta_k$ \tabularnewline
Quasiparticle & 准粒子 & 产生算符 $\alpha^\dagger$ \tabularnewline
Product Wave Function & 直积波函数 & 附录 E，独立粒子波函数 \tabularnewline
Slater Determinant & Slater 行列式 & 费米子多体波函数 \tabularnewline
Density Matrix & 密度矩阵 & $\rho_{kk'} = \langle c_{k'}^\dagger c_k \rangle$，附录 D \tabularnewline
Pairing Tensor & 配对张量 & $\kappa_{kk'} = \langle c_{k'} c_k \rangle$ \tabularnewline
Generalized Density Matrix & 广义密度矩阵 & $\mathcal{R} = \begin{bmatrix} \rho & \kappa \\ -\kappa^* & 1-\rho^* \end{bmatrix}$ \tabularnewline
Bogoliubov Transformation & Bogoliubov 变换 & $u, v$ 振幅矩阵 \tabularnewline
Bloch-Messiah Theorem & Bloch-Messiah 定理 & 附录 E，将广义 Bogoliubov 变换分解为典范基 \tabularnewline
Canonical Basis & 典范基 / 标准基 & 对角化密度矩阵 $\rho$ 的单粒子表象 \tabularnewline
Thouless Theorem & Thouless 定理 & 附录 E，联接两个非正交准粒子真空态 \tabularnewline
Onishi Formula & Onishi 公式 & 附录 E，直积态重叠的行列式表达式 \tabularnewline
Wick's Theorem & Wick 定理 & 算符乘积转正规序与缩并 \tabularnewline
Normal Ordering & 正规序 & $: A B C :$ \tabularnewline
Contraction & 缩并 & 费米子线缩并 \tabularnewline
Particle-Hole (ph) Channel & 粒子-空穴通道 & $\text{ph}$ 相互作用，激发算符 $a_m^\dagger a_i$ \tabularnewline
Particle-Particle (pp) Channel & 粒子-粒子通道 & $\text{pp}$ 相互作用，配对相互作用 \tabularnewline
Hole-Hole (hh) Channel & 空穴-空穴通道 & $\text{hh}$ 相互作用 \tabularnewline
Green's Function & Green 函数 & 附录 F，单体与两体传播子 \tabularnewline
\end{longtable}

\section*{3. 核结构、形变与转动运动 (Structure, Deformation \& Rotation)}

\begin{longtable}{p{4.8cm} p{4.8cm} p{5.0cm}}
\toprule
\textbf{英文原名 (En)} & \textbf{规范中文译名 (Zh)} & \textbf{常见符号 / 备注} \tabularnewline
\midrule
\endfirsthead
\toprule
\textbf{英文原名 (En)} & \textbf{规范中文译名 (Zh)} & \textbf{常见符号 / 备注} \tabularnewline
\midrule
\endhead
\bottomrule
\endfoot
\bottomrule
\endlastfoot
Magic Numbers & 幻数 & 2, 8, 20, 28, 50, 82, 126 \tabularnewline
Nuclear Deformation & 核形变 & 多极展开参数 $\alpha_{\lambda\mu}$ \tabularnewline
Quadrupole / Octupole Deformation & 四极形变 / 八极形变 & $\lambda=2$ / $\lambda=3$ \tabularnewline
Hexadecapole Deformation & 十六极形变 & $\lambda=4$ \tabularnewline
Hill-Wheeler Coordinates & Hill-Wheeler 坐标 & $\beta, \gamma$ 形状参数 \tabularnewline
Prolate Deformation & 长椭球形变 & $\gamma = 0^\circ$ \tabularnewline
Oblate Deformation & 扁椭球形变 & $\gamma = 60^\circ$ \tabularnewline
Triaxial Deformation & 三轴形变 & $0^\circ < \gamma < 60^\circ$ \tabularnewline
Moment of Inertia & 转动惯量 & $\mathcal{J}$，刚体转动惯量 / 流体转动惯量 / Cranking 转动惯量 \tabularnewline
Rigid-Body Moment of Inertia & 刚体转动惯量 & $\mathcal{J}_{\text{rig}}$ \tabularnewline
Cranking Model & Cranking 模型 / 推引模型 & 描述转动核的 Inglis 理论，第 3、7 章 \tabularnewline
Rotating Frame & 转动参考系 / 体固坐标系 & 坐标系与原子核形状共同转动 \tabularnewline
Coriolis Force & Coriolis 力 & $-\boldsymbol{\omega} \cdot \boldsymbol{j}$ \tabularnewline
Coriolis Anti-Pairing & CAP 效应 / Coriolis 反配对效应 & 强转动下配对能隙塌缩 \tabularnewline
Backbending & 回弯现象 & 高自旋态下转动惯量剧增的现象 \tabularnewline
Rotational Alignment & 转动顺排 & 高 $j$ 轨道核子角动量向转动轴顺排 \tabularnewline
Particle-Rotor Model & 粒子加转子模型 & 第 3 章 \tabularnewline
Band Crossing & 能带交叉 & 基态带 (g-band) 与顺排带 (s-band) 交叉 \tabularnewline
Laboratory System & 实验室坐标系 & 空间固定坐标系 \tabularnewline
Body-Fixed System & 体固坐标系 / 内秉坐标系 & 附录 A \tabularnewline
Euler Angles & Euler 角 & $\Omega = (\alpha, \beta, \gamma)$ \tabularnewline
\end{longtable}

\section*{4. 核力、核物质与相互作用 (Nuclear Forces \& Interactions)}

\begin{longtable}{p{4.8cm} p{4.8cm} p{5.0cm}}
\toprule
\textbf{英文原名 (En)} & \textbf{规范中文译名 (Zh)} & \textbf{常见符号 / 备注} \tabularnewline
\midrule
\endfirsthead
\toprule
\textbf{英文原名 (En)} & \textbf{规范中文译名 (Zh)} & \textbf{常见符号 / 备注} \tabularnewline
\midrule
\endhead
\bottomrule
\endfoot
\bottomrule
\endlastfoot
Nucleon-Nucleon Interaction & 核子-核子相互作用 & NN 相互作用 \tabularnewline
Bare Force / Realistic Force & 裸核力 / 现实核力 & Hamada-Johnston, Reid, Bonn 势 \tabularnewline
Effective Interaction & 有效相互作用 & 平均场中使用的有效力，如 Skyrme, Gogny \tabularnewline
Repulsive Core / Hard Core & 排斥核心 / 硬芯 & 短程强排斥势 ($r \lesssim 0.5\,\text{fm}$) \tabularnewline
Soft Core & 软芯 & 连续平滑可微的短程排斥 \tabularnewline
One-Boson-Exchange Potential & 单玻色子交换势 & OBEP，基于介子场论 \tabularnewline
Yukawa Potential & Yukawa 势 & 介子交换势 \tabularnewline
Tensor Force & 张量力 & $S_{12} = 3(\boldsymbol{\sigma}_1 \cdot \hat{\boldsymbol{r}})(\boldsymbol{\sigma}_2 \cdot \hat{\boldsymbol{r}}) - \boldsymbol{\sigma}_1 \cdot \boldsymbol{\sigma}_2$ \tabularnewline
Spin-Orbit Force & 自旋-轨道力 & $\boldsymbol{l} \cdot \boldsymbol{s}$ \tabularnewline
Isospin & 同位旋 & $\boldsymbol{\tau}$，质子与中子统一定义 \tabularnewline
Charge Independence & 电荷无关性 & 同位旋标量性质 \tabularnewline
Charge Symmetry & 电荷对称性 & $pp$ 与 $nn$ 相互作用的等价性 \tabularnewline
Nuclear Matter & 核物质 & 无穷大均匀费米流体 \tabularnewline
Saturation Density & 饱和密度 & $\rho_0 \simeq 0.16 \sim 0.17\,\text{fm}^{-3}$ \tabularnewline
Incompressibility & 不可压缩模量 / 压缩系数 & $K \simeq 200 \sim 250\,\text{MeV}$ \tabularnewline
Symmetry Energy & 对称能 & 结合能中与中子过剩 $(N-Z)^2/A$ 相关的项 \tabularnewline
Surface Energy & 表面能 & $a_S A^{2/3}$ \tabularnewline
Coulomb Energy & 库仑能 & 质子间静电排斥能 \tabularnewline
Bethe-Weizsäcker Formula & Bethe-Weizsäcker 公式 / 半经验质量公式 & 第 1 章 \tabularnewline
G-Matrix & G 矩阵 & Brueckner 理论处理硬芯的重整化相互作用 \tabularnewline
\end{longtable}

\section*{5. 对关联与超流核术语 (Pairing \& Superfluidity)}

\begin{longtable}{p{4.8cm} p{4.8cm} p{5.0cm}}
\toprule
\textbf{英文原名 (En)} & \textbf{规范中文译名 (Zh)} & \textbf{常见符号 / 备注} \tabularnewline
\midrule
\endfirsthead
\toprule
\textbf{英文原名 (En)} & \textbf{规范中文译名 (Zh)} & \textbf{常见符号 / 备注} \tabularnewline
\midrule
\endhead
\bottomrule
\endfoot
\bottomrule
\endlastfoot
Cooper Pair & Cooper 对 & 时间反演态对 $(k, \bar{k})$ \tabularnewline
Pairing Force & 配对力 / 对力 & 常数单极对力 $-G \sum a_k^\dagger a_{\bar{k}}^\dagger a_{\bar{k}'} a_{k'}$ \tabularnewline
Seniority Quantum Number & 辛弱数 & $v$（未配对核子数） \tabularnewline
Quasi-spin Operators & 准自旋算符 & $S_+, S_-, S_0$ \tabularnewline
Occupation Probability & 占据几率 & $v_k^2$ \tabularnewline
Vacancy Probability & 空额几率 & $u_k^2 = 1 - v_k^2$ \tabularnewline
Pairing Gap & 配对能隙 & $\Delta$ \tabularnewline
Gap Equation & 能隙方程 & $\Delta = G \sum_k u_k v_k$ \tabularnewline
Chemical Potential & 化学势 / Fermi 能 & $\lambda$（约束平均粒子数 $\langle \hat{N} \rangle = A$） \tabularnewline
Quasiparticle Energy & 准粒子能量 & $E_k = \sqrt{(\epsilon_k - \lambda)^2 + \Delta^2}$ \tabularnewline
Even-Odd Mass Staggering & 奇偶质量差 / 奇偶交错 & 配对效应导致的基态结合能系统性差异 \tabularnewline
Pairing Vibration & 对振动 & 跨闭壳两核子转移激发的模式 \tabularnewline
\end{longtable}

\section*{6. 集体激发与动力学 (Excitation, RPA, TDHF, Semiclassical)}

\begin{longtable}{p{4.8cm} p{4.8cm} p{5.0cm}}
\toprule
\textbf{英文原名 (En)} & \textbf{规范中文译名 (Zh)} & \textbf{常见符号 / 备注} \tabularnewline
\midrule
\endfirsthead
\toprule
\textbf{英文原名 (En)} & \textbf{规范中文译名 (Zh)} & \textbf{常见符号 / 备注} \tabularnewline
\midrule
\endhead
\bottomrule
\endfoot
\bottomrule
\endlastfoot
Giant Resonance & 巨共振 & GR（核子集体相干运动激发） \tabularnewline
Giant Dipole Resonance & 巨偶极共振 & GDR \tabularnewline
Giant Quadrupole Resonance & 巨四极共振 & GQR \tabularnewline
Giant Monopole Resonance & 巨单极共振（呼吸模） & GMR (Breathing Mode) \tabularnewline
Isoscalar Mode & 同位旋标量模 & 质子与中子同相位运动 ($\Delta T = 0$) \tabularnewline
Isovector Mode & 同位旋矢量模 & 质子与中子反相位运动 ($\Delta T = 1$) \tabularnewline
Forward Amplitude & 正向振幅 & $X_{mi}$ (粒子-空穴激发振幅) \tabularnewline
Backward Amplitude & 反向振幅 & $Y_{mi}$ (地面关联 / 粒子-空穴去激发振幅) \tabularnewline
Spurious State / Spurious Mode & 赝态 / 赝模（虚假态） & 质心运动或平移/转动对称性自发破缺引入的无物理能量激发 \tabularnewline
Goldstone Mode & Goldstone 模 / 零能模 & 恢复破缺对称性对应的零激发能本征模 \tabularnewline
Energy-Weighted Sum Rule & 能量加权求和规则 & EWSR \tabularnewline
Boson Expansion Method & 玻色展开方法 & 第 9 章 \tabularnewline
Marumori Expansion & Marumori 展开 & 第 9 章 \tabularnewline
Belyaev-Zelevinsky Expansion & Belyaev-Zelevinsky 展开 & 第 9 章 \tabularnewline
Generator Coordinate & 生成坐标 & $q$ (形变参数、对能隙等动力学自由度) \tabularnewline
Hill-Wheeler Equation & Hill-Wheeler 方程 & GCM 中的积分本征值方程 \tabularnewline
Projection Before Variation & 变分前投影 & PBV / VAP (Variation After Projection) \tabularnewline
Projection After Variation & 变分后投影 & PAV \tabularnewline
Particle Number Projection & 粒子数投影 & PNP \tabularnewline
Angular Momentum Projection & 角动量投影 & AMP \tabularnewline
WKB Approximation & WKB 近似 & Wentzel-Kramers-Brillouin 半经典波函数 \tabularnewline
Thomas-Fermi Model & Thomas-Fermi 模型 & 费米气体局域动量空间近似 \tabularnewline
Extended Thomas-Fermi & 扩展 Thomas-Fermi 模型 & ETF，包含 $\hbar^2, \hbar^4$ 梯度修正 \tabularnewline
Wigner Transform & Wigner 变换 & 相空间分布函数 $f(\boldsymbol{r}, \boldsymbol{p})$ \tabularnewline
Wigner Function & Wigner 函数 & 经典相空间量子对应 \tabularnewline
\end{longtable}

\section*{7. 跃迁概率、角动量代数与电磁极矩 (Transitions \& Electromagnetism)}

\begin{longtable}{p{4.8cm} p{4.8cm} p{5.0cm}}
\toprule
\textbf{英文原名 (En)} & \textbf{规范中文译名 (Zh)} & \textbf{常见符号 / 备注} \tabularnewline
\midrule
\endfirsthead
\toprule
\textbf{英文原名 (En)} & \textbf{规范中文译名 (Zh)} & \textbf{常见符号 / 备注} \tabularnewline
\midrule
\endhead
\bottomrule
\endfoot
\bottomrule
\endlastfoot
Reduced Transition Probability & 约化跃迁几率 / 约化跃迁概率 & $B(E\lambda), B(M\lambda)$ \tabularnewline
Wigner-Eckart Theorem & Wigner-Eckart 定理 & 附录 B，分解矩阵元为几何因子与动力学约化矩阵元 \tabularnewline
Electric Multipole Operator & 电多极算符 & $\mathcal{M}(E\lambda, \mu)$ \tabularnewline
Magnetic Multipole Operator & 磁多极算符 & $\mathcal{M}(M\lambda, \mu)$ \tabularnewline
Electric Quadrupole Moment & 电四极矩 & $Q_0$ (内秉电四极矩), $Q_{\text{spec}}$ (谱学电四极矩) \tabularnewline
Magnetic Dipole Moment & 磁偶极矩 & $\mu$ \tabularnewline
Gyromagnetic Ratio & 旋磁比 / $g$ 因子 & $g_l, g_s, g_R$ \tabularnewline
Weisskopf Single-Particle Unit & Weisskopf 单粒子单位 & W.u.，单粒子跃迁强度的基准单位 \tabularnewline
Selection Rules & 选择定则 & 宇称、角动量三角定则 \tabularnewline
\end{longtable}

\section*{8. 经典文献与人名惯用列表 (Proper Names \& Literature)}

\begin{multicols}{2}

\subsection*{原著作者Peter Ring, Peter Schuck}

\subsection*{文献著作}
\begin{itemize}
\item Bohr \& Mottelson {[}BM 75{]}
\item Eisenberg \& Greiner {[}EG 70{]}
\item de Shalit \& Feshbach {[}SF 74{]}
\item Edmonds {[}Ed 57{]}
\end{itemize}

\subsection*{理论物理人名}
\begin{itemize}
\item Hartree, Fock
\item Bogoliubov
\item Slater
\item Bloch, Messiah
\item Thouless
\item Onishi
\item Wick
\item Bardeen, Cooper, Schrieffer
\item Strutinsky
\item Nilsson
\item Inglis
\item Skyrme, Gogny
\item Brueckner, Bethe, Goldstone
\item Yukawa
\item Coriolis
\item Clebsch, Gordan
\item Wigner, Eckart
\item Weisskopf
\item Hill, Wheeler
\item Marumori, Belyaev, Zelevinsky
\item Wentzel, Kramers, Brillouin (WKB)
\item Thomas, Fermi
\item Pauli
\item Dirac
\item Schrödinger
\item Hamilton
\item Feynman
\end{itemize}
\end{multicols}
